% GENERATED FILE. Edit canonical lessons, then run npm run generate:books.
% canonical-source-hash: fnv1a64:a4204f8011914a01
% canonical-lessons: SA-C281-ekam-yanapatram-dadatu, SA-C281-agrimah, SA-C281-mama-syutah-nastah, SA-C281-ausadham, SA-S228-letter-au, SA-R281-at-the-station

\chapter{At the Station, and When Things Go Wrong}
\label{ch:sa-at-the-station-and-when-things-go-wrong}
\hlchaptermodality{281}

\begin{chapteropening}
\textbf{By the end of this chapter:} \emph{I can ask for tickets and when the next train leaves (ekaṃ yānapatraṃ dadātu; agrimaṃ dhūmayānaṃ kadā?), and say what went wrong (mama syūtaḥ naṣṭaḥ; dhūmayānaṃ gatam; ahaṃ rugṇaḥ asmi).}

Complete SA-R281-at-the-station: buy a ticket, ask when the train leaves, and say what went wrong.
\end{chapteropening}

\section[\texorpdfstring{\sk{एकं} \sk{यानपत्रं} \sk{ददातु} (ekaṃ yānapatraṃ dadātu)}{एकं यानपत्रं ददातु (ekaṃ yānapatraṃ dadātu)}]{\sk{एकं} \sk{यानपत्रं} \sk{ददातु} (ekaṃ yānapatraṃ dadātu) — a ticket, please}

\label{lesson:SA-C281-ekam-yanapatram-dadatu}



\begin{quote}
Before the new one: say \textquotedblleft{}where is the station?\textquotedblright{} (\emph{sthānakaṃ kutra asti?}) Then say \textquotedblleft{}a ticket\textquotedblright{}. (\emph{yānapatram}.) You are at the counter now.
\end{quote}

\subsection*{You'll want to know: tickets}
\begin{quote}
\textbf{\sk{एकं} \sk{यानपत्रं} \sk{ददातु।}}
\end{quote}

\begin{quote}
\emph{ekaṃ yānapatraṃ dadātu.} — \textquotedblleft{}One ticket, please.\textquotedblright{}
\end{quote}

\textbf{\sk{एकम्}} (\emph{ekam}), \textquotedblleft{}one\textquotedblright{}, agrees with \textbf{\sk{यानपत्रम्}}, a neuter noun. For two, the number is \textbf{\sk{द्वे}} (\emph{dve}) and the noun takes the dual, \textbf{\sk{यानपत्रे}} (\emph{yānapatre}):

\begin{quote}
\textbf{\sk{कृपया} \sk{द्वे} \sk{यानपत्रे} \sk{ददातु।}}
\end{quote}

\begin{quote}
\emph{kṛpayā dve yānapatre dadātu.} — \textquotedblleft{}Two tickets, please.\textquotedblright{}
\end{quote}

\begin{quote}
\textbf{\sk{मूल्यं} \sk{कियत्}?}
\end{quote}

\begin{quote}
\emph{mūlyaṃ kiyat?} — \textquotedblleft{}How much is it?\textquotedblright{}
\end{quote}

\subsection*{You'll want to know: when does it leave?}
The future asks when: \textbf{\sk{गमिष्यति}} (\emph{gamiṣyati}), \textquotedblleft{}will go\textquotedblright{}, is \textquotedblleft{}will leave\textquotedblright{}.

\begin{quote}
\textbf{\sk{धूमयानं} \sk{कदा} \sk{गमिष्यति}?}
\end{quote}

\begin{quote}
\emph{dhūmayānaṃ kadā gamiṣyati?} — \textquotedblleft{}When will the train leave?\textquotedblright{}
\end{quote}

\begin{quote}
\textbf{\sk{सायं} \sk{गमिष्यति।}}
\end{quote}

\begin{quote}
\emph{sāyaṃ gamiṣyati.} — \textquotedblleft{}It will leave in the evening.\textquotedblright{}
\end{quote}

\begin{quote}
\textbf{\sk{लोकयानं} \sk{कुत्र} \sk{तिष्ठति}?}
\end{quote}

\begin{quote}
\emph{lokayānaṃ kutra tiṣṭhati?} — \textquotedblleft{}Where does the bus stop?\textquotedblright{}
\end{quote}

\subsection*{Your turn}
Say these aloud:

\begin{itemize}
\raggedright
  \item \textquotedblleft{}The train will leave in the morning\textquotedblright{} — \emph{dhūmayānaṃ prātaḥ gamiṣyati}
  \item \textquotedblleft{}The bus will leave in the morning\textquotedblright{} — \emph{lokayānaṃ prātaḥ gamiṣyati}
  \item \textquotedblleft{}The plane will leave in the morning\textquotedblright{} — \emph{vimānaṃ prātaḥ gamiṣyati}
\end{itemize}

\begin{tcolorbox}[breakable,colback=teal!4,colframe=teal!35!black,title={Before you move on}]
\textquotedblleft{}One ticket, please\textquotedblright{}? (\textbf{\sk{एकं} \sk{यानपत्रं} \sk{ददातु}}.) \textquotedblleft{}Two tickets\textquotedblright{}? (\textbf{\sk{द्वे} \sk{यानपत्रे}}.) \textquotedblleft{}When will the train leave?\textquotedblright{} (\textbf{\sk{धूमयानं} \sk{कदा} \sk{गमिष्यति}?})
\end{tcolorbox}
\section[\texorpdfstring{\sk{अग्रिमः} (agrimaḥ)}{अग्रिमः (agrimaḥ)}]{\sk{अग्रिमः} (agrimaḥ) — next}

\label{lesson:SA-C281-agrimah}



\begin{quote}
Before the new one: say \textquotedblleft{}when will the train leave?\textquotedblright{} (\emph{dhūmayānaṃ kadā gamiṣyati?}) That one has gone. You want the next.
\end{quote}

\subsection*{You'll want to know: \sk{अग्रिमः}}
\textbf{\sk{अग्रिमः}} — \emph{agrimaḥ} — \textquotedblleft{}next\textquotedblright{}, the one coming up. It agrees with its noun: \textbf{\sk{अग्रिमम्}} (\emph{agrimam}) for a neuter noun such as \textbf{\sk{धूमयानम्}}.

\begin{quote}
\textbf{\sk{अग्रिमं} \sk{धूमयानं} \sk{कदा}?}
\end{quote}

\begin{quote}
\emph{agrimaṃ dhūmayānaṃ kadā?} — \textquotedblleft{}When is the next train?\textquotedblright{}
\end{quote}

\textbf{\sk{अग्रिमः} \sk{सप्ताहः}} (\emph{agrimaḥ saptāhaḥ}) is \textquotedblleft{}next week\textquotedblright{}.

\begin{quote}
\textbf{\sk{अहम्} \sk{अग्रिमे} \sk{सप्ताहे} \sk{गमिष्यामि।}}
\end{quote}

\begin{quote}
\emph{aham agrime saptāhe gamiṣyāmi.} — \textquotedblleft{}I will go next week.\textquotedblright{}
\end{quote}

\textbf{\sk{अग्रिमे} \sk{सप्ताहे}} (\emph{agrime saptāhe}) is \textquotedblleft{}in the next week\textquotedblright{}: \emph{-e} says \textquotedblleft{}in\textquotedblright{}, for a time as for a place.

\begin{cousinweb}
\textbf{\sk{अग्रिमः}} (\emph{agrimaḥ}) is built on \textbf{\sk{अग्र}} (\emph{agra}), \textquotedblleft{}the foremost point, the front\textquotedblright{}: Monier-Williams gives \emph{agrima} as \textquotedblleft{}being in front, preceding\textquotedblright{}. The next train is the one in front. CDIAL 92 lists Punjabi \emph{aggiõ}, \textquotedblleft{}first\textquotedblright{}.
\end{cousinweb}

\subsection*{Your turn}
Say these aloud:

\begin{itemize}
\raggedright
  \item \emph{agrimaḥ}
  \item \textquotedblleft{}When is the next train?\textquotedblright{} — \emph{agrimaṃ dhūmayānaṃ kadā?}
  \item \textquotedblleft{}When is the next bus?\textquotedblright{} — \emph{agrimaṃ lokayānaṃ kadā?}
\end{itemize}

\begin{tcolorbox}[breakable,colback=teal!4,colframe=teal!35!black,title={Before you move on}]
What does \textbf{\sk{अग्रिमः}} mean? (\textquotedblleft{}Next\textquotedblright{}.) \textquotedblleft{}The next train\textquotedblright{}? (\textbf{\sk{अग्रिमं} \sk{धूमयानम्}}.)
\end{tcolorbox}
\section[\texorpdfstring{\sk{मम} \sk{स्यूतः} \sk{नष्टः} (mama syūtaḥ naṣṭaḥ)}{मम स्यूतः नष्टः (mama syūtaḥ naṣṭaḥ)}]{\sk{मम} \sk{स्यूतः} \sk{नष्टः} (mama syūtaḥ naṣṭaḥ) — when things go wrong}

\label{lesson:SA-C281-mama-syutah-nastah}



\begin{quote}
Say \textquotedblleft{}the next train\textquotedblright{}. (\emph{agrimaṃ dhūmayānam}.) Then say \textquotedblleft{}is lost\textquotedblright{}. (\emph{naśyati}.) Now something has gone wrong.
\end{quote}

\subsection*{You'll want to know: lost}
\begin{quote}
\textbf{\sk{मम} \sk{स्यूतः} \sk{नष्टः।}}
\end{quote}

\begin{quote}
\emph{mama syūtaḥ naṣṭaḥ.} — \textquotedblleft{}My bag is lost.\textquotedblright{}
\end{quote}

\textbf{\sk{नष्टः}} (\emph{naṣṭaḥ}), \textquotedblleft{}lost\textquotedblright{}, is the \emph{-ta} form of \textbf{\sk{नश्यति}} (\emph{naśyati}), \textquotedblleft{}is lost\textquotedblright{}. It describes, so it agrees with the thing: \textbf{\sk{स्यूतः}} is masculine, so \textbf{\sk{नष्टः}}.

\begin{quote}
\textbf{\sk{मम} \sk{यानपत्रं} \sk{नष्टम्।}}
\end{quote}

\begin{quote}
\emph{mama yānapatraṃ naṣṭam.} — \textquotedblleft{}My ticket is lost.\textquotedblright{}
\end{quote}

\begin{quote}
\textbf{\sk{मम} \sk{माला} \sk{नष्टा।}}
\end{quote}

\begin{quote}
\emph{mama mālā naṣṭā.} — \textquotedblleft{}My garland is lost.\textquotedblright{}
\end{quote}

\subsection*{You'll want to know: missed, and ill}
\begin{quote}
\textbf{\sk{धूमयानं} \sk{गतम्।}}
\end{quote}

\begin{quote}
\emph{dhūmayānaṃ gatam.} — \textquotedblleft{}The train has gone.\textquotedblright{}
\end{quote}

\textbf{\sk{गतम्}} (\emph{gatam}) is \textbf{\sk{गतः}} for a neuter noun: the train is gone.

\begin{quote}
\textbf{\sk{अहं} \sk{रुग्णः} \sk{अस्मि।}}
\end{quote}

\begin{quote}
\emph{ahaṃ rugṇaḥ asmi.} — \textquotedblleft{}I am ill.\textquotedblright{}
\end{quote}

A woman says \textbf{\sk{अहं} \sk{रुग्णा} \sk{अस्मि}} (\emph{ahaṃ rugṇā asmi}). And to ask for help, the polite \textbf{-\sk{तु}} on \textquotedblleft{}does\textquotedblright{}:

\begin{quote}
\textbf{\sk{कृपया} \sk{साहाय्यं} \sk{करोतु।}}
\end{quote}

\begin{quote}
\emph{kṛpayā sāhāyyaṃ karotu.} — \textquotedblleft{}Please help.\textquotedblright{}
\end{quote}

\subsection*{Your turn}
Say these aloud:

\begin{itemize}
\raggedright
  \item \textquotedblleft{}My bag is lost\textquotedblright{} — \emph{mama syūtaḥ naṣṭaḥ}
  \item \textquotedblleft{}My ticket is lost\textquotedblright{} — \emph{mama yānapatraṃ naṣṭam}
  \item \textquotedblleft{}My book is lost\textquotedblright{} — \emph{mama pustakaṃ naṣṭam}
\end{itemize}

\begin{tcolorbox}[breakable,colback=teal!4,colframe=teal!35!black,title={Before you move on}]
\textquotedblleft{}My bag is lost\textquotedblright{}? (\textbf{\sk{मम} \sk{स्यूतः} \sk{नष्टः}}.) Why \textbf{\sk{नष्टम्}} with \textbf{\sk{यानपत्रम्}}? (It is neuter.) \textquotedblleft{}The train has gone\textquotedblright{}? (\textbf{\sk{धूमयानं} \sk{गतम्}}.)
\end{tcolorbox}
\section[\texorpdfstring{\sk{औषधम्} (auṣadham)}{औषधम् (auṣadham)}]{\sk{औषधम्} (auṣadham) — a medicine}

\label{lesson:SA-C281-ausadham}



\begin{quote}
Before the new one: say \textquotedblleft{}I am ill\textquotedblright{}. (\emph{ahaṃ rugṇaḥ asmi}.) Then \textquotedblleft{}you should go to the doctor\textquotedblright{}, and \textquotedblleft{}the doctor will give advice\textquotedblright{}. (\emph{bhavān vaidyaṃ gacchet}; \emph{vaidyaḥ upadeśaṃ dāsyati}.)
\end{quote}

\subsection*{You'll want to know: \sk{औषधम्}}
\textbf{\sk{औषधम्}} — \emph{auṣadham} — \textquotedblleft{}a medicine\textquotedblright{}. It is a neuter noun.

\begin{quote}
\textbf{\sk{औषधं} \sk{कुत्र} \sk{अस्ति}?}
\end{quote}

\begin{quote}
\emph{auṣadhaṃ kutra asti?} — \textquotedblleft{}Where is the medicine?\textquotedblright{}
\end{quote}

\begin{quote}
\textbf{\sk{वैद्यः} \sk{औषधं} \sk{दास्यति।}}
\end{quote}

\begin{quote}
\emph{vaidyaḥ auṣadhaṃ dāsyati.} — \textquotedblleft{}The doctor will give medicine.\textquotedblright{}
\end{quote}

\begin{quote}
\textbf{\sk{भवान्} \sk{औषधं} \sk{पिबेत्।}}
\end{quote}

\begin{quote}
\emph{bhavān auṣadhaṃ pibet.} — \textquotedblleft{}You should take the medicine.\textquotedblright{}
\end{quote}

\textbf{\sk{पिबेत्}} (\emph{pibet}) is \textbf{\sk{पिबति}} (\emph{pibati}), \textquotedblleft{}drinks\textquotedblright{}, with the \emph{-et} of advice.

\begin{cousinweb}
\textbf{\sk{औषधम्}} (\emph{auṣadham}) is built on \textbf{\sk{ओषधि}} (\emph{oṣadhi}), \textquotedblleft{}a herb\textquotedblright{}: Monier-Williams gives \emph{auṣadha} as \textquotedblleft{}consisting of herbs\textquotedblright{}, and then herbs used as medicine. It lives on as Oriya \emph{usuā}, \textquotedblleft{}medicine\textquotedblright{} (CDIAL 2571).
\end{cousinweb}

\subsection*{Your turn}
Say these aloud:

\begin{itemize}
\raggedright
  \item \emph{auṣadham}
  \item \textquotedblleft{}I am ill. Where is the medicine?\textquotedblright{} — \emph{ahaṃ rugṇaḥ asmi. auṣadhaṃ kutra asti?}
  \item \textquotedblleft{}You should take the medicine\textquotedblright{} — \emph{bhavān auṣadhaṃ pibet}
\end{itemize}

\begin{tcolorbox}[breakable,colback=teal!4,colframe=teal!35!black,title={Before you move on}]
What does \textbf{\sk{औषधम्}} mean? (\textquotedblleft{}A medicine\textquotedblright{}.) \textquotedblleft{}Where is the medicine?\textquotedblright{} (\textbf{\sk{औषधं} \sk{कुत्र} \sk{अस्ति}?})
\end{tcolorbox}
\section[\texorpdfstring{\sk{औ} (au)}{औ (au)}]{\sk{औ} — the vowel of \sk{ौ}, standing alone}

\label{lesson:SA-S228-letter-au}



\begin{quote}
Before the new one: say \textbf{\sk{औषधम्}} (\emph{auṣadham}), \textquotedblleft{}a medicine\textquotedblright{}. Then name the sign \textbf{\sk{ौ}} and the letter \textbf{\sk{अ}}. (\emph{au}, \emph{a}.)

The sign \textbf{\sk{ौ}} has a standing form, and \textbf{\sk{औषधम्}} opens with it.
\end{quote}

\subsection*{Script you'll notice: \sk{औ}}
\textbf{\sk{औ}} — \emph{au}.

It is an \textbf{independent vowel}: the full-size letter \emph{au} wears when it opens a word. On a consonant the same vowel is the sign \textbf{\sk{ौ}}, as in \textbf{\sk{द्वौ}} (\emph{dvau}).

What it is made of:

\begin{itemize}
\raggedright
  \item the joined upper and lower bowls of \sk{अ}
  \item a middle shoulder running into the inner stem
  \item the inner stem
  \item the trailing stem
  \item two separate upper arcs
  \item the top shirorekhā
\end{itemize}

You have met it here:

\begin{itemize}
\raggedright
  \item \textbf{\sk{औषधम्}} \emph{auṣadham} — a medicine
\end{itemize}

\subsection*{Writing: \sk{औ}}
\hlblockfigure{\detokenize{figures/SA-S228-letter-au-filmstrip.pdf}}{How \sk{औ} is written, stroke by stroke}

\begin{itemize}
\raggedright
  \item \textbf{1.} curve right around the upper bowl and continue down and around the lower bowl without lifting
  \item \textbf{2.} lift and sweep the middle shoulder right
  \item \textbf{3.} without lifting, climb up the inner stem to the headline
  \item \textbf{4.} descend the inner stem
  \item \textbf{5.} lift and descend the trailing stem
  \item \textbf{6.} lift at the trailing stem's headline junction and sweep the lower upper arc upward and left
  \item \textbf{7.} lift at the same junction and sweep the taller upper arc upward and left
  \item \textbf{8.} lift and draw the top shirorekhā left-to-right
\end{itemize}

\textbf{Pen lifts: 5.} The pen comes up 5 times and no more.

\begin{quote}
The source shows seven movements; native writers draw \sk{औ} as six pen-down strokes (60\% of HP Labs India's native-writer samples, 49 of 82). This ductus follows the seven movements in order and direction, lifting five times: the middle shoulder runs on, climbing the inner stem and descending it, and every other movement starts with a lift.
\end{quote}

\begin{quote}
This is one attested teaching order and not a national standard — handwriting here is taught with school-to-school variation. Source: Saurmandal, ‘Devanagari \sk{औ} stroke order.svg’, panels 1–7, Wikimedia Commons, 5 August 2023.
\end{quote}

\subsection*{Your turn}
\begin{itemize}
\raggedright
  \item \emph{Look:} at this, and find \sk{औ} in it
\end{itemize}

\begin{quote}
\sk{औषधम्}
\end{quote}

\begin{itemize}
\raggedright
  \item \emph{Trace it:} \sk{औ} three times, saying \emph{au} as you finish each one
  \item \emph{Look:} at \textbf{\sk{औषधं} \sk{कुत्र} \sk{अस्ति}?} and find \sk{औ} once more
\end{itemize}

\begin{tcolorbox}[breakable,colback=teal!4,colframe=teal!35!black,title={Before you move on}]
Which character is this — \sk{औ}? What sound does it carry? (\emph{au}.) Its form on a consonant? (\textbf{\sk{ौ}}.) Which word did it come from? (\textbf{\sk{औषधम्}}, \textquotedblleft{}a medicine\textquotedblright{}.)
\end{tcolorbox}
\section[(review)]{Review — at the station, and when things go wrong}

\label{lesson:SA-R281-at-the-station}



\begin{quote}
Say \textquotedblleft{}where is the station?\textquotedblright{} and \textquotedblleft{}go straight on\textquotedblright{}. (\emph{sthānakaṃ kutra asti?}, \emph{ṛju gacchatu}.)
\end{quote}

\subsection*{Your turn}
\begin{itemize}
\raggedright
  \item \emph{Say it:} One ticket, please. How much is it?
  \item \emph{Say it:} When is the next train? It will leave in the evening.
  \item \emph{Say it:} My bag is lost, and the train has gone.
  \item \emph{Say it:} I am ill. Where is the medicine?
  \item \emph{Say it:} Please help.
  \item \emph{Read it:} \textbf{\sk{औषधम्}} aloud, and name the sign its first vowel takes on a consonant
\end{itemize}

\begin{tcolorbox}[breakable,colback=teal!4,colframe=teal!35!black,title={Before you move on}]
\textbf{\sk{द्वे} \sk{यानपत्रे}}: how many tickets? (Two.) \textbf{\sk{मम} \sk{माला} \sk{नष्टा}}: why \textbf{\sk{नष्टा}}? (\textbf{\sk{माला}} is feminine.) Now the other way round: \textquotedblleft{}the next bus\textquotedblright{}? (\textbf{\sk{अग्रिमं} \sk{लोकयानम्}}.)
\end{tcolorbox}
